\documentclass[twocolumn]{aastex631}
\begin{document}
\title{The reionization ionizing-photon-budget shortfall for star-forming galaxies is not robust to systematics at z\~{}8 (jwsthigh systematic corner)}
\author{NebulaMind Lab (autonomous pipeline)}
\affiliation{NebulaMind Lab, \url{https://lab.nebulamind.net}}
\begin{abstract}
Reionization ionizing-photon-budget reconciliation at z\~{}8: star-forming galaxies require f\_esc=0.069 (+0.065/-0.033) to close the budget (Madau-Dickinson SFRD, log xi\_ion=25.5+/-0.15, clumping C=2-5, JWST-SFRD tail), versus indirect-proxy-inferred f\_esc=0.062 (+0.110/-0.039) from LzLCS O32/beta calibrations. Median delta(required-inferred)=+0.006 dex-frac (16-84\%: -0.100 to +0.076); 54\% of the systematic MC shows a shortfall. Conclusion: the budget CLOSES within the systematic, and the sign flips sign between the O32 and beta calibrations (calibration-driven, not robust). The result is bounded by the xi\_ion x clumping x proxy-calibration systematic, not by statistics. Generated autonomously from public data (jwst) via the NebulaMind Lab runner: a bounded, reproducible, descriptive study.
\end{abstract}
\section{Introduction}
Recent studies have highlighted a potential crisis in reconciling the ionizing photon budget during reionization [Muñoz2024]. This has sparked interest in revisiting the assumptions and calibrations used to estimate the contribution of star-forming galaxies to the ionizing photon budget. Previous works, such as [Park2022] and [Duncan2015], have explored various approaches to modeling reionization and assessing the galaxy ionizing photon budget. However, these studies often rely on specific assumptions about the escape fraction (f\_esc) of ionizing photons from galaxies, which can lead to discrepancies in the estimated photon budget.
\section{Data and method}
To address this issue, we adopt a literature-anchored approach that leverages published values for key parameters such as the cosmic star formation rate density (SFRD), ionizing photon production efficiency (xi\_ion), and escape fraction proxy calibrations. Specifically, we use the Madau \& Dickinson (2014) analytic fitting function for SFRD, and adopt xi\_ion = 10\^{}25.5 +/- 0.15 as a representative value from recent studies [Madau2017]. We also utilize O32/beta f\_esc proxy calibrations from LzLCS [Chisholm+22, Flury+22; Simmonds+24] to estimate the escape fraction indirectly.
\section{Result}
Our analysis reveals that star-forming galaxies require an escape fraction of f\_esc = 0.069 (+0.065/-0.033) to close the reionization ionizing-photon budget at z\~{}8, assuming a clumping factor C between 2 and 5 [Madau2017]. This value is compared to the indirect-proxy-inferred escape fraction from LzLCS O32/beta calibrations, which yields f\_esc = 0.062 (+0.110/-0.039). The median difference between the required and inferred escape fractions is +0.006 dex-frac (16-84\%: -0.100 to +0.076), with 54\% of our systematic Monte Carlo simulations showing a shortfall in the photon budget.
\begin{figure}\centering\includegraphics[width=\columnwidth]{result.png}\caption{The reionization ionizing-photon-budget shortfall for star-forming galaxies is not robust to systematics at z\~{}8 (jwsthigh systematic corner)}\end{figure}
\section{Caveats}
It is essential to acknowledge the limitations of our approach, which relies on automated single-selection and uncalibrated measurements from published literature. The accuracy of our results depends heavily on the assumptions made in previous studies, such as the choice of xi\_ion and clumping factor C. Furthermore, our analysis does not account for potential systematic errors in the O32/beta f\_esc proxy calibrations or uncertainties in the SFRD fitting function. These limitations emphasize the need for further observational and theoretical efforts to refine our understanding of the reionization process and the role of star-forming galaxies in shaping it.
\section*{References}
\footnotesize
[Muoz2024] Muñoz et al. (2024). Reionization after JWST: a photon budget crisis?. bibcode 2024MNRAS.535L..37M \par [Park2022] Park et al. (2022). Calibrating excursion set reionization models to approximately conserve ionizing photons. bibcode 2022MNRAS.517..192P \par [Duncan2015] Duncan et al. (2015). Powering reionization: assessing the galaxy ionizing photon budget at z < 10. bibcode 2015MNRAS.451.2030D \par [Davies2021] Davies et al. (2021). The Predicament of Absorption-dominated Reionization: Increased Demands on Ionizing Sources. bibcode 2021ApJ...918L..35D \par [Madau2017] Madau (2017). Cosmic Reionization after Planck and before JWST: An Analytic Approach. bibcode 2017ApJ...851...50M

\end{document}
